% [原文 p.20]
\subsection{部分拓扑扭曲}

我们的目标是把 6d $(2,0)$ 理论放到 $\R^4\times C_2$ 上，其中 $C_2$ 是任意的带穿刺点黎曼面。太天真地这样做不会保持 $\R^4$ 的庞加莱对称性之外的任何对称性。我们来解释一个程序------\emph{部分拓扑扭曲}------它使得无论 $C_2$ 如何都能保持 4d $\cN=2$ 超对称。

\para{拓扑扭曲的一般讨论。}首先我们从几个观点一般地评述超对称理论的拓扑扭曲 [166]。

当把场论放到弯曲背景上时，度规 $g_{\mu\nu}$ 充当能动张量 $T^{\mu\nu}$ 的源。对超对称场论，$T^{\mu\nu}$ 通常与超对称流 $S_\mu^{\ \alpha}$ 及 R 对称流 $J^\mu$ 属于同一个多重态。它们也可以与源 $\psi_\mu^{\ \alpha}$ 与 $A_\mu$ 耦合。部分拓扑扭曲就是令 $S_\mu^{\ \alpha}=0$、并取 $A_\mu$ 等于由 $g_{\mu\nu}$ 导出的自旋联络。示意地说，在 $g,\psi,A$ 的某组背景值附近作线性化，当这些场改变时拉格朗日量变化为
\begin{equation}
\delta\cL = T\delta g + J\delta A = T\delta g - J\partial(\delta g) \simeq (T+\partial J)\delta g .
\tag{3.1}
\end{equation}
第二步用了我们的选取：$A$ 与度规的导数有关；最后一步作了分部积分。

% [Tutor Note: p.20 式 (3.1) 的「示意」推导补全：$A_\mu$ 取为自旋联络，它对 $g_{\mu\nu}$ 是一阶的，
%  故 $\delta A\sim\partial(\delta g)$（这正是原文第二步）。于是 $J\delta A = -J\partial(\delta g)$，
%  对 $J\partial(\delta g)$ 分部积分并把全导数项丢掉，得 $\delta\cL\simeq(T+\partial J)\delta g$，
%  即 $T^{\mu\nu}_{\text{twist}} = T^{\mu\nu}+\partial J$。核对方法：把 (3.1) 与 (3.3) 对照------
%  (3.3) 里 $\frac14(\epsilon^{\mu\rho}\partial_\rho J^\nu_{12}+\epsilon^{\nu\rho}\partial_\rho J^\mu_{12})$
%  正是 $\partial J$ 在 $x^4,x^5$ 平面上的具体实现（两个 $\epsilon$ 收缩给出 $\partial_4J_5,\partial_5J_4$ 型两项）。
%  原文用 $\simeq$ 而非 $=$，是因为指标、对称化与整体系数都被略去。]

于是拓扑扭曲相当于在把理论放到非平凡背景度规之前，把能动张量从 $T$ 重新定义为 $T_{\mathrm{twist}}=T+\partial J$。这个扭曲把能动张量 $T^{\mu\nu}$ 与 R 对称流 $J^\mu$ 混在一起，但值得记住：它不影响理论在平坦空间中的任何可观测量，只影响我们所谓的能动张量。通过改变能动张量，它改变了理论被放到弯曲空间上的方式。

能动张量的一个职责是记录庞加莱对称性：$T^{\mu\nu}$ 是平移对称性的守恒流，而 $x^{[\mu}T^{\nu]\rho}$ 是旋转的守恒流。由于扭曲把 $T$ 平移一个全导数，它只是平移对称流的一次改进变换（improvement transformation），并不改变相应的守恒荷------动量算符。相反，它对所谓的旋转有非平凡效应：扭曲的旋转作用为一次旋转加上一次 R 对称变换，因为示意地有 $xT\mapsto xT+x\partial J=xT-J+\partial(xJ)$，而 $\partial(xJ)$ 是一次改进变换。尽管如此，平移与扭曲旋转之间的对易子仍与标准庞加莱代数中的一致。

超荷会怎样？它们通常在原来的旋转下、以及在 R 对称变换下按旋量变换。在对角地嵌入的新旋转下，超荷通常会分裂为一个标量超荷 $Q$ 与一个矢量。由于 $Q$ 是标量，能动张量是 $Q$-闭的，于是用扭曲的能动张量把理论放到弯曲流形上会保持超对称 $Q$。下一步通常是限制到 $Q$-上同调中的算子。在许多情形下扭曲的能动张量是 $Q$-恰当的，即 $T^{\mu\nu}_{\mathrm{twist}}=\{Q,G^{\mu\nu}\}$，其中 $G^{\mu\nu}$ 是某个超对称生成元；于是它在 $Q$-上同调中消失，关联函数由

% [原文 p.21]
拓扑量子场论（TQFT）描述。对我们所考虑的扭曲，这不会发生，反而会留下非平凡的局域动力学。

\para{6d 理论的部分拓扑扭曲。}我们所用的\emph{部分}拓扑扭曲只把\emph{某些} R 对称性混入\emph{某些}旋转对称性。为定义我们所用的具体扭曲，考虑分别保持乘积 $\R^{1,3}\times\R^2$ 两个因子的旋转 $\fso(1,3)\times\fso(2)_{\mathrm{old}}$，并考虑 R 对称性 $\fso(5)_{\mathrm R}$ 的块对角子代数 $\fso(2)_{\mathrm R}\times\fso(3)_{\mathrm R}\subset\fso(5)_{\mathrm R}$。我们把扭曲旋转定义成对角地嵌入 $\fso(2)_{\mathrm{old}}\times\fso(2)_{\mathrm R}$，即把下列对称性当作我们（扭曲后的）洛伦兹与 R 对称性：
\begin{equation}
\fso(1,3)\times\fso(2)_{\mathrm{twist}}\times\fso(3)_{\mathrm R} .
\tag{3.2}
\end{equation}
这是通过把能动张量改为
\begin{equation}
T^{\mu\nu}_{\mathrm{twist}} = T^{\mu\nu}_{\mathrm{old}} + \frac{1}{4}\left(\epsilon^{\mu\rho}\partial_\rho J^{\nu}_{12} + \epsilon^{\nu\rho}\partial_\rho J^{\mu}_{12}\right) ,
\tag{3.3}
\end{equation}
来实现的，其中 $J_{12}$ 是 $\fso(2)_{\mathrm R}$ 的 R 对称旋转生成元，而 $\epsilon^{\mu\nu}=\delta^\mu_4\delta^\nu_5-\delta^\nu_5\delta^\mu_4$ 是 $\R^2$ 因子上的 Levi-Civita 张量。

% [Tutor Note: p.21 式 (3.3) 之后 $\epsilon^{\mu\nu}$ 的定义在原文里印作
%  $\epsilon^{\mu\nu}=\delta^\mu_4\delta^\nu_5-\delta^\nu_5\delta^\mu_4$，两项完全相同，故该式恒等于零------
%  这是原文的笔误，正确形式应为 $\epsilon^{\mu\nu}=\delta^\mu_4\delta^\nu_5-\delta^\mu_5\delta^\nu_4$。
%  判断依据有二：(i) 相减的两项必须不同，否则 $\R^2$ 上的 Levi-Civita 张量恒为零，与 (3.3) 和习题 3.1
%  都要用到它矛盾；(ii) 直接核对分量：正确形式给出 $\epsilon^{45}=\delta^4_4\delta^5_5-\delta^4_5\delta^5_4=1$、
%  $\epsilon^{54}=\delta^5_4\delta^4_5-\delta^5_5\delta^4_4=-1$，即 $\epsilon^{45}=-\epsilon^{54}=1$，
%  这正是二维 Levi-Civita 的要求。译文正文按 STYLE_GUIDE §9.1 照抄原文的错误形式，正确形式见本条；
%  Lead 若希望正文直接排正确形式，只需把那句里的 $\delta^\nu_5\delta^\mu_4$ 换成 $\delta^\mu_5\delta^\nu_4$。]

\begin{exercise}{3.1}
验证 (3.3) 把 $x^4,x^5$ 旋转流 $x^{[4}T^{5]\mu}$ 平移了 $J_{12}$（至多差全导数，即一项改进项），从而使扭曲旋转是旋转与 R 对称性的组合。
\end{exercise}

让我们追踪超对称在扭曲、然后再紧化的过程中如何变化。在 $\R^{1,5}$ 的 $\fso(1,5)$ 旋转与 $\fso(5)_{\mathrm R}$ R 对称性下，庞加莱超对称按各自的旋量表示变换，记作 $(\mathbf{4},\mathbf{4})$，并带一个我们为简单起见隐去的辛实条件。每个 6d Weyl 旋量，即 $\fso(1,5)$ 的每个表示 $\mathbf{4}$，在 $\fso(1,3)$ 下分解为一对相反手征的 4d Weyl 旋量 $(\mathbf{2},\mathbf{1})\oplus(\mathbf{1},\mathbf{2})$，而这些旋量在 $\fso(2)_{\mathrm{old}}$ 下的电荷相反，为 $1/2$ 与 $-1/2$。$\fso(5)$ 的每个旋量 $\mathbf{4}$ 分解为 $\fso(3)_{\mathrm R}$ 的两个 $\mathbf{2}$，其 $\fso(2)_{\mathrm R}$ 电荷为 $\pm1/2$。合起来我们如下记之，下标表示在两个 $\fso(2)$ 代数下的电荷：
\begin{equation}
\begin{aligned}
(\mathbf{4},\mathbf{4}) &= \Big((\mathbf{2},\mathbf{1})_{\tfrac12}\oplus(\mathbf{1},\mathbf{2})_{-\tfrac12}\Big)\otimes\Big(\mathbf{2}_{\tfrac12}\oplus\mathbf{2}_{-\tfrac12}\Big)\\
&= (\mathbf{2},\mathbf{1},\mathbf{2})_{\tfrac12,\tfrac12}\oplus(\mathbf{2},\mathbf{1},\mathbf{2})_{\tfrac12,-\tfrac12}\oplus(\mathbf{1},\mathbf{2},\mathbf{2})_{-\tfrac12,\tfrac12}\oplus(\mathbf{1},\mathbf{2},\mathbf{2})_{-\tfrac12,-\tfrac12} .
\end{aligned}
\tag{3.4}
\end{equation}
由构造，$\fso(2)_{\mathrm{twist}}$ 下的电荷是 $\fso(2)_{\mathrm{old}}$ 与 $\fso(2)_{\mathrm R}$ 下电荷之和。于是，在 $\R^{1,5}$ 的 $\fso(1,3)\times\fso(3)_{\mathrm R}\times\fso(2)_{\mathrm{twist}}$ 对称性下，庞加莱超荷变换为
\begin{equation}
(\mathbf{2},\mathbf{1};\mathbf{2})_1\oplus(\mathbf{2},\mathbf{1};\mathbf{2})_0\oplus(\mathbf{1},\mathbf{2};\mathbf{2})_0\oplus(\mathbf{1},\mathbf{2};\mathbf{2})_{-1} .
\tag{3.5}
\end{equation}
我们分别记它们为
\begin{equation}
Q^{\alpha A}_z ,\ Q^{\alpha A} ,\ \bar Q^{\dot\alpha A} ,\ \bar Q^{\dot\alpha A}_{\bar z} ,
\tag{3.6}
\end{equation}

% [Tutor Note: p.21 式 (3.4)--(3.6) 的下标对应关系可自查：$\fso(2)_{\mathrm{old}}$ 电荷 $+1/2$ 与
%  $\fso(2)_{\mathrm R}$ 电荷 $+1/2$ 之和为 $+1$，正好给出 (3.5) 的第一项 $(\mathbf{2},\mathbf{1};\mathbf{2})_1$，
%  即 $Q_z$；而 $-\tfrac12+\tfrac12=0$ 给出中间两个标量超荷 $Q,\bar Q$。核对方法：把 (3.4) 的第二行四项
%  按 twist 电荷（两个下标之和）分组，得到 $+1,0,0,-1$，与 (3.5) 的四项一一对应，符号也一致。]

% [原文 p.22]
其中 $\alpha,\dot\alpha,A$ 取值 $1,2$，分别是 $\fso(1,3)$ 两种手征的旋量指标与 $\fso(3)_{\mathrm R}$ 的旋量指标；而 $z$ 是 $\R^2$ 因子上的复坐标，用来记录第一个与最后一个超荷 $Q_z,\bar Q_{\bar z}$ 的 $\fso(2)_{\mathrm{twist}}$ 电荷 $\pm1$。

中间两个超荷 $Q,\bar Q$ 在 $\fso(2)_{\mathrm{twist}}$ 旋转下是标量，因此把 $\R^2$ 上的度规形变到任何弯曲度规都保持这些超荷。合起来，在部分拓扑扭曲下紧化到 $\R^{1,3}\times C$ 上，我们得到一个保持 $\mathfrak{iso}(1,3)$ 庞加莱对称性、超荷 $Q^{\alpha A}$ 与 $\bar Q^{\dot\alpha A}$、以及 $\fso(3)_{\mathrm R}=\fsu(2)$ R 对称性的系统。它们共同构成 4d $\cN=2$ 庞加莱超对称代数。

在 $C$ 的尺寸为零的极限下，我们一般地得到 4d $\cN=2$ 理论。\footnote{有限面积的 $C$ 上的系统有一定的真空模空间；在把面积趋于零的标度极限中，必须指定在哪个真空附近展开。如果 $C$ 有“足够多”的柄或穿刺点，则其 Higgs 分支有一个最大对称点，在这个点附近展开是自然的，于是 4d $\cN=2$ 极限良定义。如果 $C$ 是“穿刺点太少”的球面、或没有穿刺点的环面，则不存在最大对称点，情形更微妙，如 [8] 所解释。}扭曲 (3.3) 不保持 $T$ 的无迹性，因此尽管原来的 6d 旋转对称性可以延拓为共形对称性，扭曲后的旋转对称性却不然。在零面积极限下，4d 共形对称性可以恢复，于是我们得到一个 SCFT------除非 $C$ 的穿刺点处的数据带有内禀质量标度。

\subsection{Coulomb 分支}

4d $\cN=2$ 理论的 Coulomb 分支由给 Coulomb 分支算子------即 4d 理论中被所有反手征庞加莱超荷 $\bar Q^{\dot\alpha A}$ 湮灭的（规范不变）算子------一个 VEV 来描述。让我们从 6d 理论 $\cX(\fg)$ 出发来确认这些算子，大致遵循 [167, 第 3 节]。

重要的是，由 (3.13) 得到的 Coulomb 分支 $\cB$ \emph{只依赖于 $C$ 的复结构}，而不依赖其度规。这让我们可以用各种方式形变黎曼面以理解所得的 4d 理论，它也是 AGT 对应中 $C$ 上出现 2d CFT 对象的基础。

一般超对称考虑允许 4d $\cN=2$ 理论 [168] 的真空模空间是混合分支 $\cC_\alpha\times\cH_\alpha$ 的并，其中包括纯 Coulomb 分支与纯 Higgs 分支作为特例。特殊 Kähler 流形 $\cC_\alpha$ 由 Coulomb 分支算子参数化，超 Kähler 流形 $\cH_\alpha$ 由 Higgs 分支算子参数化。class S 理论的 Higgs 分支手征环关系已在 [169--172] 中探讨。像 [173, 174] 那样确定 class S 理论的所有分支一般是困难的，故我们只关注 Coulomb 分支（此时 $\cH_\alpha$ 是一个点）。

\para{Coulomb 分支算子。}$\cX(\fg)$ 的真空由标量场 $\Phi_I$（$I=6,\dots,10$）在 $\fg$ 的 Cartan 子代数中的 VEV（模掉 Weyl 群）参数化。给定真空中的低能理论由这些场的涨落、以及旋量 $\lambda$ 与自对偶二形式 $B$ 描述。在我们所关心的、恰在扭曲之前的对称性代数 $\big(\fso(1,3)\times\fso(3)_{\mathrm R}\big)\times\fso(2)_{\mathrm{old}}\times\fso(2)_{\mathrm R}$ 下，这些场变换为

% [原文 p.23]
\begin{equation}
\begin{aligned}
B &\in (\mathbf{3},\mathbf{1},\mathbf{1})_{0,0}\oplus(\mathbf{1},\mathbf{3},\mathbf{1})_{0,0}\oplus(\mathbf{2},\mathbf{2},\mathbf{1})_{\pm1,0}\oplus(\mathbf{1},\mathbf{1},\mathbf{1})_{0,0},\\
\lambda &\in (\mathbf{2},\mathbf{1},\mathbf{2})_{\tfrac12,\pm\tfrac12}\oplus(\mathbf{1},\mathbf{2},\mathbf{2})_{-\tfrac12,\pm\tfrac12},
&&\Phi_z := \Phi_6+i\Phi_7\in(\mathbf{1},\mathbf{1},\mathbf{1})_{0,1},\\
\Phi_8,\Phi_9,\Phi_{10} &\in(\mathbf{1},\mathbf{1},\mathbf{3})_{0,0},
&&\Phi_{\bar z} := \Phi_6-i\Phi_7\in(\mathbf{1},\mathbf{1},\mathbf{1})_{0,-1} .
\end{aligned}
\tag{3.7}
\end{equation}

另一方面，超荷 $\bar Q^{\dot\alpha A}$ 按 $(\mathbf{1},\mathbf{2},\mathbf{2})_{-1/2,1/2}$ 变换。我们断言
\begin{equation}
\bar Q^{\dot\alpha A}\Phi_z = 0
\tag{3.8}
\end{equation}
因为 $\lambda$ 的任何分量都没有合适的 $\fso(2)_{\mathrm R}$ 电荷 $3/2$。

其实，我们应该使用相应的规范不变算子，例如经典情形 $\fsu(N)$ 与 $\fso(2N)$ 中的迹 $\Tr(\Phi^j_z)$。这些就是 $\fg$ 的 Casimir 不变量，即各种阶数 $d_k$（$k=1,\dots,\rk\fg$）的多项式 $P_k(\Phi_z)$。具体地，对经典规范群，这些被 $\bar Q^{\dot\alpha A}$ 湮灭的规范不变算子是
\begin{equation}
\begin{gathered}
\Tr(\Phi^j_z),\qquad j=2,3,4,\dots,N \qquad\qquad \text{for } \fsu(N),\\
\Tr(\Phi^j_z),\qquad j=2,4,6,\dots,2N-2,\ \text{and}\quad \operatorname{Pfaff}(\Phi_z)\qquad \text{for } \fso(2N).
\end{gathered}
\tag{3.9}
\end{equation}
（我们回忆 Pfaffian 是行列式的平方根。）供参考，$\fsu(N)$ 的 Casimir 阶数是 $2,3,\dots,N$；$\fso(2N)$ 的是 $2,4,6,\dots,2N-2$ 与 $N$；$\mathfrak{e}_6$ 的是 $2,5,6,8,9,12$；$\mathfrak{e}_7$ 的是 $2,6,8,10,12,14,18$；$\mathfrak{e}_8$ 的是 $2,8,12,14,18,20,24,30$。

% [Tutor Note: p.23 的 Casimir 阶数表可逐项核对（判据：阶数个数必须等于该代数的秩）：
%  $\fsu(N)$ 秩 $N-1$，阶数 $2,\dots,N$ 共 $N-1$ 个 ✓；$\fso(2N)$ 秩 $N$，阶数 $2,4,\dots,2N-2$ 与 $N$
%  共 $(N-1)+1=N$ 个 ✓（$N$ 偶时会出现两个同为 $N$ 的 Casimir，正是 p.23 脚注 26 所说的退化情形）；
%  $\mathfrak{e}_6,\mathfrak{e}_7,\mathfrak{e}_8$ 秩分别为 $6,7,8$，所给列表长度分别为 $6,7,8$ ✓。
%  另一个独立锚点：$\mathfrak{e}_8$ 的最大阶数 $30$ 等于其 Coxeter 数 $h=30$，$\mathfrak{e}_7$ 的最大阶数 $18=h$，
%  $\mathfrak{e}_6$ 的最大阶数 $12=h$。]

经常方便的是把 $\Phi_z$ 换成 $\Phi_z\dd z$，以吸收 $z$ 指标、得到一个张量。于是我们在全纯曲线（又称黎曼面）$C$ 上使用 $d_k$ 阶微分 $P_k(\Phi_z)\dd z^{d_k}$。一个稍有不同的基更实用：例如对 $\fsu(N)$ 展开
\begin{equation}
\det(X-\Phi_z\dd z) = X^N - \sum_{j=2}^{N}\cO_j X^{N-j} .
\tag{3.10}
\end{equation}

\begin{exercise}{3.2}
验证 $\cO_2=\Tr(\Phi^2_z/2)\dd z^2$、$\cO_3=\Tr(\Phi^3_z/3)\dd z^3$，并（如果愿意）验证 $\cO_4=\Tr(\Phi^4_z/4)\dd z^4-\cO_2^2/2$。为什么没有 $\cO_1$？
\end{exercise}

% [Tutor Note: p.23 习题 3.2 的两问都可直接验算。(i) 为什么没有 $\cO_1$：$\fsu(N)$ 的
%  $\Phi_z$ 无迹，$\Tr\Phi_z=0$，故定义 $\cO_1$ 的那个系数恒为零（也可理解为 $\fsu(N)$ 的 u(1) 已解耦，
%  (3.10) 的求和从 $j=2$ 起）。(ii) $\cO_4$ 的表达式即 Newton 恒等式：设 $p_k=\Tr\Phi^k$、$e_k$ 为
%  特征值的初等对称多项式，$k e_k=\sum_{i=1}^{k}(-1)^{i-1}e_{k-i}p_i$。取 $p_1=0$ 得
%  $e_2=-p_2/2$、$e_3=p_3/3$、$4e_4=e_3p_1-e_2p_2+e_1p_3-e_0p_4=p_2^2/2-p_4$，
%  故 $e_4=p_4/4-p_2^2/8=\Tr(\Phi^4/4)-(\Tr(\Phi^2/2))^2/2=\Tr(\Phi^4/4)-\cO_2^2/2$ ✓，
%  与原文一致。]

\para{Coulomb 分支。}我们可以给 $\cO_j$ 什么 VEV 来定义一个真空？记它为\footnote{对 $\fso(4K)$ 情形这个记号略有不适定，因为此时有两个同阶数 $2K$ 的 Casimir，从而有两个 $2K$ 阶微分：由 $\Phi_z$ 的幂的迹定义的 $\phi_{2K}$，以及 $\tilde\phi_{2K}=\langle\operatorname{Pfaff}(\Phi_z)\rangle\dd z^{2K}$。}
\begin{equation}
\phi_j := \langle\cO_j\rangle .
\tag{3.11}
\end{equation}
它应当沿 $\R^{1,3}$ 为常数，以免破坏庞加莱对称性。接着我们用反交换子 $\{\bar Q^{\dot\alpha A},\bar Q^{\dot\beta B}_{\bar z}\}\sim\epsilon^{\dot\alpha\dot\beta}\epsilon^{AB}\partial_{\bar z}$ 推出 $\phi_j$ 必须全纯地依赖于 $z$：
\begin{equation}
\epsilon^{\dot\alpha\dot\beta}\epsilon^{AB}\partial_{\bar z}\phi_j \sim \Big\langle \bar Q^{\dot\alpha A}\big(\bar Q^{\dot\beta B}_{\bar z}\cO_j\big)\Big\rangle + \Big\langle \bar Q^{\dot\beta B}_{\bar z}\big(\bar Q^{\dot\alpha A}\cO_j\big)\Big\rangle = 0 .
\tag{3.12}
\end{equation}

% [原文 p.24]
第一项消失，因为扭曲紧化保持超荷 $\bar Q$。第二项由 $\cO_j$ 的构造而消失。

于是 4d $\cN=2$ 理论的 Coulomb 分支由 $C$ 上阶数为 $d_k$ 的微分 $\phi_{d_k}$（$k=1,\dots,\rk\fg$）参数化。当没有穿刺点时，
\begin{equation}
\cB = \bigoplus_{k=1}^{r} H^0(C, K^{\otimes d_k}) ,\qquad \phi_j\in H^0(C, K^{\otimes j}) ,
\tag{3.13}
\end{equation}
其中 $K$ 是曲线 $C$ 上的典范丛，$H^0(C,\cL)$ 是 $C$ 上线丛 $\cL$ 的截面组成的向量空间。当 $C$ 有穿刺点时，截面 $\phi_{d_k}$ 在每个穿刺点处有指定的行为。从第 4 节开始，针对给出通常 4d $\cN=2$ 规范理论的具体 $C$ 选择，我们会解释如何把参数化 (3.13) 与用 4d $\cN=2$ 矢量多重态中的标量给出的通常描述联系起来。

\para{Hitchin 系统。}我们想直观地说，4d Coulomb 分支由伴随值全纯一形式 $\Phi_z\dd z$（一个假想中的 $H^0(C,K\otimes\fg)$ 元素）的“VEV”模掉规范变换来参数化。当然，非规范不变算子的 VEV 没有意义（或者按你的观点自动为零），所以谈论它们是一种语言滥用。然而在我们的情形存在所谓 Hitchin 场（或 Higgs 场）的构造：一个全纯一形式 $\varphi=\varphi_z\dd z$，其分量 $\varphi_z\in\fg$，其 Casimir 在 $\fsu(N)$ 情形给出 $\Tr(\Phi^j_z)\dd z^j$，在其它情形类似。为方便起见，在某些解释中我们偶尔用 $\varphi$ 而不是规范不变量 $\phi_j$。

故事还要长一点：人们把 4d 理论进一步紧化在 $S^1$ 上。3d 理论的 Coulomb 分支真空由 $C$ 上 Hitchin 系统的解 $(A,\varphi)$ 给出，
\begin{equation}
F+[\varphi,\bar\varphi]=0 ,\qquad \bar\partial_A\varphi=0 ,\qquad \partial_A\bar\varphi=0 ,
\tag{3.14}
\end{equation}
模掉 $G$ 规范变换。所得的 Coulomb 分支 $\cM$（Hitchin 模空间）通过把 $(A,\varphi)$ 映到 $\varphi$ 的 Casimir，给出到 4d 理论 Coulomb 分支 $\cB$ 的投影。Hitchin 方程 (3.14) 等价于 $G_\C$ 联络 $A+\varphi_z\dd z+\bar\varphi_{\bar z}\dd\bar z$ 的平坦性加上一个规范固定条件，因此 $\cM$ 也可以描述为 $C$ 上复 $G_\C$ 平坦联络模掉 $G_\C$ 规范变换所得的模空间。

% [Tutor Note: p.24 式 (3.14) 的三条 Hitchin 方程中，第二、三条互为厄米共轭（$\bar\varphi=\varphi^\dagger$、
%  $\partial_A$ 与 $\bar\partial_A$ 共轭），故独立的只有前两条：$\bar\partial_A\varphi=0$（全纯性）与
%  $F+[\varphi,\varphi^\dagger]=0$（"实"条件）。原文把它们并列写出是为了保持 $A\leftrightarrow\bar A$ 的对称性。
%  另注意 $F$ 的具体系数/符号依规范化而定（原文未给系数，只写 $+$），习惯上取
%  $F_{z\bar z}+[\varphi_z,\varphi_{\bar z}]=0$。核对方法：把第一式按 $z,\bar z$ 分量写出，
%  并用 Hitchin 方程可推出 $A+\varphi\,\dd z+\bar\varphi\,\dd\bar z$ 平坦，这与原文
%  「等价于 $G_\C$ 联络的平坦性 + 规范固定」的说法一致。]

% [Tutor Note: p.24 两处可留意：(i) 原文 "modulo $G_\C$ gauge transformaions" 拼写有误（缺 t），
%  应为 "transformations"；译文照常译出。(ii) 同一页 $G_\C$ 联络那一句里，原文把 $\Phi_z\dd z$ 的
%  微分排成直立体 $\mathrm{d}$、而把 $+\bar\varphi_{\bar z}d\bar z$ 的微分排成斜体 $d$，前后不一致；
%  按 STYLE_GUIDE 统一用 $\dd$。另外 (3.13) 的求和指标用 $k$（$\phi_{d_k}$），而同一式右半却用 $\phi_j$、
%  原文未加说明，$j$ 与 $k$ 指的是同一组阶数（$j=d_k$ 遍历 Casimir 阶数），读者不必当成两组指标。]

\para{关于红外的行为。}6d 理论在一般真空中的低能极限由取值于真空模空间的阿贝尔 6d $(2,0)$ 理论给出。同样地，在由 (3.13) 中一组给定微分 $\phi_j$ 所描述的 Coulomb 分支真空中，4d $\cN=2$ 理论的有效描述包含描述 $\phi_j$ 涨落的无质量标量场。与 $B_{\mu\nu}$ 和 $\lambda$ 的类似维数约化一起，这些标量场组成 4d $\cN=2$ 阿贝尔矢量多重态。

有多少个？标量场以 Coulomb 分支 $\cB$ 为目标，故我们应期待红外描述是一个规范群为 $\mathrm{U}(1)^{\dim_\C\cB}$ 的 4d $\cN=2$ 规范理论。在 Coulomb 分支的特定点处还有额外的无质量粒子，它们带这个规范群的荷。

% 边界说明：上句起于 p.24、收尾于 p.25 第一行，按 BRIEF「跨页句子开始者写完」由本块写完。
% p.25 从下一句 "The Coulomb branch typically features points that are even more singular…" 起
% 归 tr-02（含 (3.15) 及其后全部内容），本块不写。
